% !PRINTED 10
\weberpage{10}
%<<CONT
$A''$ enthalten sind, nämlich die Brüche von der Form $a+a'$ und
$b+b'$. Es ist ferner $\alpha''$ grösser als $\alpha$ und $\alpha'$.
Denn $A''$ enthält zunächst alle $A$, und wenn $a'$ ein beliebiger
Bruch in $A'$ ist, so kann man $a$ in $A$ so wählen, dass das Element
$a+a'$ von $A''$ in $B$ enthalten ist; also ist $A''$ umfassender als
$A$. Die durch rationale Brüche $\mu,\mu'$ erzeugten Schnitte ergeben
durch Addition den durch $\mu+\mu'$ erzeugten Schnitt.

\begin{center}
\rule{0.18\linewidth}{0.4pt}
\end{center}

Wir gehen nun über zu der Definition der \gesperrt{Verhältnisse}, die
von Alters her als Grundlage der Zahlenlehre betrachtet werden, und
folgen dabei zunächst Euklid.\setcounter{footnote}{0}\footnote{Elemente, Buch V.}

Wenn man die Elemente einer messbaren Menge $\mathfrak M$ zu Paaren
verbindet, und diese Paare an sich als Elemente betrachtet, so
entsteht eine neue Mannigfaltigkeit; wir bezeichnen ein solches Paar
mit $a:b$, oder auch mit $a/b$, unterscheiden aber, wenn $a$ und $b$
verschiedene Elemente sind, $a:b$ von $b:a$, und nennen $a$ den
\gesperrt{Zähler} und $b$ den \gesperrt{Nenner} von $a:b$. Diese
Paare wollen wir Verhältnisse nennen und wollen diese neue Menge nun
ordnen und messbar machen.

Nehmen wir zunächst an, dass zwei ganze Zahlen $m,n$ existiren, so
dass $na=mb$ wird, wie es z.~B. immer der Fall ist, wenn $\mathfrak M$
das System der natürlichen Zahlen ist, oder wenn $a,b$ zwei
commensurable Strecken sind; dann ist, wenn $p,q$ zwei andere ganze
Zahlen sind, dann und nur dann $qa=pb$, wenn $mq=np$ ist. Das
Zahlenpaar $p,q$ ist durch diese Forderung vollständig bestimmt, wenn
noch die Bedingung hinzukommt, dass $p,q$ relative Primzahlen sein
sollen. Dann kann, wenn $h$ eine beliebige ganze Zahl ist, $m=hp$,
$n=hq$ sein. In diesem Falle nennen wir das Verhältniss $a:b$ ein
\gesperrt{rationales} und setzen es gleich dem rationalen Bruch $m:n$
oder $p:q$. Diese rationalen Brüche können hiernach als Verhältnisse
ganzer Zahlen aufgefasst werden.

Alle unter einander gleichen rationalen Verhältnisse bilden eine
\gesperrt{rationale Zahl}, und die rationalen Zahlen bilden, wie die
rationalen Brüche, eine geordnete, dichte und messbare
Mannigfaltigkeit.
